% ECONOMICS_PROBLEM: {"id": "L51-495", "title": "Privacy as a competition variable", "jel_code": "L51", "source_id": "OP-0076"}
% !TeX program = xelatex
\documentclass[11pt,a4paper]{article}
\usepackage[margin=23mm]{geometry}
\usepackage{fontspec}
\setmainfont{Latin Modern Roman}
\usepackage{amsmath,amssymb,mathtools}
\usepackage{xcolor,enumitem,booktabs,longtable,array,xurl,hyperref}
\definecolor{muted}{HTML}{53636C}
\hypersetup{colorlinks=true,urlcolor=blue,linkcolor=blue}
\setlength{\parindent}{0pt}
\setlength{\parskip}{5pt}
\newcommand{\R}{\mathbb R}
\newcommand{\E}{\mathbb E}
\newcommand{\N}{\mathbb N}
\newcommand{\ind}{\mathbf 1}
\newcommand{\Var}{\operatorname{Var}}
\newcommand{\Cov}{\operatorname{Cov}}
\newcommand{\supp}{\operatorname{supp}}
\DeclareMathOperator*{\argmin}{arg\,min}
\DeclareMathOperator*{\argmax}{arg\,max}
\newcommand{\entrymeta}[3]{{\sffamily\small\color{muted}\textbf{\texttt{#1}}\quad|\quad #2\par #3\par}}
\newcommand{\Problemref}[1]{\texttt{#1}}
\newcommand{\JELref}[2]{\texttt{#2} (JEL \texttt{#1})}
\begin{document}
\section*{Privacy as a competition variable}
\textbf{Problem ID:} \texttt{L51-495}\quad \textbf{JEL:} \texttt{L51}\par
% BEGIN ECONOMICS STATEMENT
\label{problem:OP-0076}
\label{atlas:prob:I6}

\noindent\textbf{Original identifier:} I6 (atlas item 76). \textbf{Classification:} Privacy and competition equilibrium research program. \textbf{Original tractability label:} Medium (editorial, not a complexity result).

\subsection*{Definitions and assumptions}

A privacy regime $r$ specifies permissible collection, retention, sharing, and targeting uses of data. Consumers have types $\theta$ determining product demand and utility from data practices. Firms choose privacy attributes, prices, advertising technology, and entry, facing regime-dependent compliance and information costs. Let $e_\sigma(r)$ be a selected equilibrium of this declared game. A welfare evaluator $w$ defines the treatment of privacy harms, informed consent, advertising nuisance, and public externalities. Consumer surplus must be based on a utility model incorporating relevant privacy effects; advertising revenue is not itself an additional social benefit after transfers are accounted for.

\subsection*{Formal research question}

For regimes $r_1,r_0$, identify or bound
\[
\Delta CS^\sigma=CS(e_\sigma(r_1),r_1)-CS(e_\sigma(r_0),r_0),\quad
\Delta W^{w,\sigma}=W^w(e_\sigma(r_1),r_1)-W^w(e_\sigma(r_0),r_0),
\]
alongside ad prices, entrant survival, and market shares. Characterize estimable conditions under which privacy gains coexist with increased concentration, and whether total welfare remains positive. If a regulatory shock changes several constraints simultaneously, identify the bundled intervention or supply exclusions separating its privacy, compliance-cost, and access-to-data channels. Report selection and valuation sensitivity where point identification fails.

\subsection*{Interpretation and scope}

Stronger privacy is a multidimensional policy description and needs a concrete regime before a counterfactual is defined. Restricting third-party tracking may affect entrants differently from firms with first-party data, while giving users control may alter product quality or participation. Consumer willingness to pay is informative under assumptions about knowledge, choice costs, and risk; observed acceptance of tracking does not reveal informed preferences automatically. Entry effects require a baseline-defined cohort or a dynamic entry model, avoiding conditioning on firms that survive the policy. Concentration can rise even when every consumer gains privacy utility, so a concentration statistic cannot settle the welfare question. Conversely, ad-price changes alone do not determine effects on advertiser profit or consumers. A useful empirical extension would combine a well-defined policy shock with consumer demand, vendor relationships, costs, and entrant outcomes. Where the shock is common to all markets, credible comparison groups and timing assumptions require particular attention. The analysis should disclose how different privacy welfare evaluators alter conclusions.

\subsection*{Source audit and status}

[S1] verifies the 2011 journal issue despite earlier online dating; [S2] is the economics-of-privacy review. The underspecified Johnson et al. (2023) citation is contextually resolved to Johnson, Shriver, and Goldberg's GDPR vendor-concentration study [S3]. That setting does not itself identify all consumer surplus or ad-funded entry effects. The joint welfare and competition target above is an editorial extension, with application-specific current status rather than a universal certified open conjecture.

\subsection*{References}

\begin{enumerate}[label={[S\arabic*]},leftmargin=*]

\item Avi Goldfarb and Catherine Tucker (2011). \emph{Privacy Regulation and Online Advertising}. Management Science 57(1), 57--71. \url{https://doi.org/10.1287/mnsc.1100.1246}\par \textit{Source role:} Evidence on privacy regulation and advertising effectiveness. \textit{Locator:} Abstract and article metadata.

\item Alessandro Acquisti, Curtis Taylor, and Liad Wagman (2016). \emph{The Economics of Privacy}. Journal of Economic Literature 54(2), 442--492. \url{https://www.aeaweb.org/articles?id=10.1257/jel.54.2.442}\par \textit{Source role:} Survey of privacy valuations and economic mechanisms. \textit{Locator:} Abstract and article metadata.

\item Garrett Johnson, Scott K. Shriver, and Samuel G. Goldberg (2023). \emph{Privacy and Market Concentration: Intended and Unintended Consequences of the GDPR}. Management Science 69(10), 5695--5721. \url{https://doi.org/10.1287/mnsc.2023.4709}\par \textit{Source role:} Evidence on GDPR and web-technology vendor concentration. \textit{Locator:} Abstract and article metadata.

\end{enumerate}

\subsection*{Common conventions: Source mathematical conventions}

Unless an entry states otherwise, agent and firm sets are finite. State and action spaces are measurable, strategies and outcome maps are measurable, probability laws are specified, and expectations exist for the targets under discussion. Infinite-horizon discount factors lie in $(0,1)$ and discounted objectives are finite. Norms, tolerances and complexity input sizes must be specified for computational comparisons. Constants or primitives that appear locally are inputs, not empirical facts asserted by this catalogue.

\textbf{Identification.} A statistical model $m\in\mathcal M$ implies an observable law $P_m$ and target $\theta(m)$. At law $P$, its identified set is
\[
\Theta(P;\mathcal M)=\{\theta(m):m\in\mathcal M,\ P_m=P\}.
\]
Point identification holds when this set is a singleton, or equivalently when $P_{m_1}=P_{m_2}$ implies $\theta(m_1)=\theta(m_2)$ for all compatible models. Sharp bounds mean bounds attained, or approached when the boundary is not attained, by compatible models. Writing this set is a definition, not a solution: the substantive task is to establish informative restrictions and derive or estimate the set. An unrestricted-confounding model may give only trivial bounds.

\textbf{Inference.} Identification, estimability and valid uncertainty quantification are separate questions. Fix $\alpha\in(0,1)$. A confidence procedure $C_n$ over a declared law class $\mathcal P$ has asymptotic uniform coverage at level $1-\alpha$ when
\[
\liminf_{n\to\infty}\inf_{P\in\mathcal P}\Pr_P\{\theta(P)\in C_n\}\geq1-\alpha.
\]
For set-identified targets the coverage event must be defined separately, for example coverage of the full identified set. If an entry uses identification-indexed models instead of uniquely indexed laws, the same coverage convention is applied over those models. Pointwise limits do not establish uniform validity, and asymptotic validity does not guarantee finite-sample precision. Sample sizes, dependence structures and weak-identification sequences are part of each application.

\textbf{Counterfactuals.} $Y_i(a)$ is a potential outcome under a specified intervention, including its timing, implementation and versions. It is not $Y_i$ conditional on voluntarily choosing $a$. A full intervention vector permits interference; shorthand scalar treatment notation does not silently impose no interference. Assignment, selection, attrition, anticipation and measurement must be specified. A claim about an instrument's local effect is not automatically a population policy effect.

\textbf{Economic equilibrium.} A counterfactual model includes best responses, constraints, feasibility, market clearing, state transitions and expectations consistent with the stated information. When several equilibria exist, either a declared selector is an input or the target is set-valued. Comparisons across policy regimes must use the stated selection convention. Reduced-form relationships are not presumed invariant after an equilibrium-changing intervention.

\textbf{Optimization and welfare.} Optimization is over the stated feasible set. If attainment is not established, $\sup$ or $\inf$ and $\varepsilon$-optimal choices replace an asserted maximizer or minimizer. For example, with value $V=\sup_{a\in\mathcal A}W(a)<\infty$, an $\varepsilon$-optimal set is $\{a\in\mathcal A:W(a)\geq V-\varepsilon\}$ for $\varepsilon>0$. Benefits, costs, risk and health or environmental outcomes can be aggregated only using declared units and conversion weights. Interpersonal utility weights, the treatment of fiscal transfers and foreign profits, discounting, and the behavioral welfare criterion are explicit inputs. A comparative-static derivative additionally requires existence and appropriate differentiability of the selected counterfactual.

\textbf{Sources.} Local labels [S1], [S2], and so forth restart in every question. Locators identify the relevant abstract, model, section or result at the level actually checked; a bibliographic or abstract-level verification is not represented as inspection of an entire proof. Working papers are distinguished from later publications and changing versions. Source summaries are paraphrased. All URL checks and editorial formulations refer to the audit date stated above.
% END ECONOMICS STATEMENT
\end{document}
